\documentclass[french]{article}

\usepackage{babel}
\usepackage{psfig}
\usepackage{comment}
\usepackage{times}
\usepackage[T1]{fontenc}
\usepackage[french]{minitoc}
\usepackage[latin1]{inputenc}
\usepackage{moreverb}
\usepackage{latexsym}
\usepackage{amssymb}
\newenvironment{smallverbatim}%
{\endgraf\small\verbatimtab}{\endverbatimtab}

\newcommand{\fig}[2]
{
 \begin{figure*}[hbt]
  \begin{center}
   \centerline{\psfig{figure=#1.eps}}
  \caption{\label{#1}  #2}
  \end{center}
 \end{figure*}
}

\newenvironment{solution}{\nopagebreak\par\vspace*{2mm}\underline{Solution~:}\em\par}{\par\vspace*{2mm}}


\pagestyle{headings}
\markright{06/03/2006 \hfill Dauphine - M1 apprentissage \hfill\ } 

\baselineskip 6pt
\parskip 6pt
\onecolumn
\textwidth 16cm
\textheight 24cm
\hoffset=-2cm
\voffset=-2.5cm
\setlength{\parindent}{0.0cm}
\unitlength 1cm
\newtheorem{defin}{Définition}[section]
\newtheorem{exampl}{Exemple}[section]
\def\picwidth{16cm}

\excludecomment{solution}

\begin{document}
\begin{center}
\Large\bf  Examen Bases de Données\\
\large\bf Lundi 6 mars 2006
\end{center}

%%%%%%%%%%%%%%%%%%%%%%%%%%%%%%%%%%%%%%%%%%%%%%%%%%%%%%%%%%%%

\vspace*{0.5cm}

\hrule

\textbf{Documents autorisés: notes de cours et polycopié}

\smallskip

\hrule

\section{Stockage et indexation (3 pts)}

On veut stocker un fichier $F$ avec 120\,000 articles d'une taille fixe de 
100 octets par article.

Questions :
\begin{enumerate}
  \item (1 point) De combien de blocs a-t-on besoin au minimum pour
           stocker tout
      le fichier si la taille d'un bloc est 8\,192 octets, sachant
          que chaque bloc contient un entête de 150 octets
          et qu'un enregistrement ne peut pas chevaucher 2 blocs\,?


 %     Solution : 1 500 blocs

  \item (1 point) 
    On suppose maintenant que $F$ est indexé par un index non dense 
    sur \texttt{A} et
     un index dense sur un autre champ \texttt{B}. On peut
   stocker 100 entrées dans un bloc d'index et
       aucune place libre n'est laissée dans les blocs. 
 Combien d'entrées y a-t-il dans les feuilles de l'index non dense\,? 
dans les feuilles de l'index dense\,? 
\begin{solution}
Le fichier $F$ contient 1\,500 blocs et et trié sur \texttt{A}.
 L'index non dense contient donc 1\,500
 entrées, une par bloc. L'index dense contient 120\,000 entrées, une par
article.
\end{solution}

  \item (1 point) Supposons que l'index non dense a deux niveaux, racine
     comprise, et que seule la racine est en mémoire.  On se place dans le pire
     des cas où il faut une E/S pour lire une feuille d'index et une E/S pour
     lire un bloc du fichier. On cherche les enregistrements pour lesquels la
     valeur de l'attribut \texttt{A} est comprise entre P1A3 et P3G5. On
     suppose qu'il y a moins de 100 enregistrements tels que \texttt{A}
     commençe par P. Combien d E/S coûte la recherche par l'index dans le pire
     des cas\,?

\begin{solution}
La recherche utilise  l'index non dense. Celui-ci tient sur 32 blocs.
La traversée de l'index part de la racine et arrive immédiatement à la
feuille d'index dont la premiere entrée est inférieure ou égale à P123 et dont
la dernière entrée est supérieure ou égale à P1A3. On recherche dans cette
feuille l'entrée la plus grande inférieure ou égale à P1A3. Cette entrée
contient l'adresse du bloc du fichier qui contient l'article de code  P1A3. 
On accède au bloc. Les articles sont triés sur \texttt{A}. Soit tous les
articles de code compris entre P1A3 et P3G5 sont dans ce bloc soit ils sont
à cheval sur deux blocs adjacents (ils ne peuvent pas être sur plus de 2 blocs
adjacents). Dans ce dernier cas, une lecture du bloc
adjacent chainé est nécessaire. Donc la recherche coûte 2 ou 3 accès
bloc.
\end{solution}

 \item Même question avec l'index dense, pour une recherche sur l'attribut
        \texttt{B}, avec les mêmes hypothèses de sélectivité.
\end{enumerate}

\section{Index et optimisation (6 points)}

Soit les tables relationnelles suivantes (les attributs qui forment 
une clé sont en gras):

\begin{itemize}
  \item Produits  (\textbf{code}, marque, desig, descr)\\
     (code : identifiant du produit, marque : marque du produit,
         desig : désignation du produit, descr : description)

  \item PrixFour (\textbf{code}, \textbf{nom}, prix)\\
   (code : code du produit, nom : nom du fournisseur)

  \item NoteMag    (\textbf{code}, \textbf{titre}, note)
     (code : code du produit,  titre : titre du magazine,
        note : note entre 1 et 10 du produit dans le magazine)
\end{itemize}

Voici une instance de la table NoteMag.

\begin{center}
\begin{tabular}{l|c|l}
code & titre & note\\
\hline
'A345' 	& 'HIFI' &	8\\
'P123' & 'Audio Expert' & 6\\
'X254' & 'HIFI' & 7\\
'K783' & 'Son \& Audio'&	3\\
'P345' & 'HIFI' & 6\\
'P512' & 'Audio Expert' & 8\\
'L830' & 'Audio Expert' & 8\\
'M240' & ''HIFI'' & 6\\
\hline
\end{tabular}
\end{center}


La table occupe plusieurs pages disque dont chacune peut contenir au maximum 3
n-uplets.

\begin{enumerate}
  \item Construire un arbre B+ sur l'attribut \texttt{code}, 
       avec 4 entrées par page au maximum.

   \item On veut indexer l'attribut \texttt{titre}. Quel est le bon
        index à utiliser\,? Décrivez-le.
   \item Soit la requête suivante :

\begin{verbatim}
select *
from NoteMag
where code between 'A000' and 'X000';
\end{verbatim}

 Pour évaluer cette requête, on suppose que le tampon de lecture ne peut
 contenir qu'une seule page et l' index tient en mémoire. Dans ce cas, est-ce
 qu'il est préférable d'utiliser l'index ou de parcourir la table
 séquentiellement ? Pourquoi ?

\begin{solution}
 Solution : Une lecture séquentielle est préférable : la sélectivité de l'index
 est très basse, et comme il s'agit d'un index dense, on risque de lire la même
 page plusieurs fois.
\end{solution}

\end{enumerate}


On donne ci-dessous une requête SQL et le plan d'exécution fourni par Oracle :

\begin{smallverbatim}
select desig, marque, prix
  from Produits, PrixFour, NoteMag
 where Produits.code=PrixFour.code
   and Produits.code=NoteMag.code
   and note > 8;
\end{smallverbatim}

Plan d'exécution :
\begin{smallverbatim}
0 SELECT STATEMENT
  1 MERGE JOIN
    2 SORT JOIN
      3 NESTED LOOPS
        4 TABLE ACCESS FULL NOTEMAG
        5 TABLE ACCESS BY INDEX ROWID PRODUITS
          6 INDEX UNIQUE SCAN A34561
    7 SORT JOIN
      8 TABLE ACCESS FULL PRIXFOUR
\end{smallverbatim}

\begin{enumerate}
 \item Existe-t-il un index ? sur quel(s) attribut(s) de quel(s) table(s) ?
\begin{solution}
      Solution : Il existe un index sur l'attribut code de la table Produits.
\end{solution}
  \item Algorithme de jointure : Expliquer en détail le plan d'exécution (accès
  aux tables, sélections, jointure, projections) 

  \item Ajout d'index : On crée un
  index sur l'attribut \texttt{note} de la table 
     \textit{NoteMag}. Expliquez les améliorations en
  terme de plan d'exécution apportées par la création de cet index.
\begin{solution}
Solution :

Après création d'index le plan est :

\begin{smallverbatim}
Plan d'execution
--------------------------------------------------------------------------------
0 SELECT STATEMENT
  1 NESTED LOOPS
    2 NESTED LOOPS
      3 TABLE ACCESS FULL PRIXFOUR
      4 TABLE ACCESS BY INDEX ROWID PRODUITS
        5 INDEX UNIQUE SCAN PRODUITS_CODE
    6 TABLE ACCESS BY INDEX ROWID NOTEMAG
      7 INDEX RANGE SCAN NOTE_MAG
\end{smallverbatim}
\end{solution}
\end{enumerate}



\section{Concurrence (6 points)}

L'exécution suivante est reçue par un SGBD utilisé pour des applications de commerce éléctronique.

$$H : r_1[x] r_2[y] w_1[x] r_2[z] r_3[z] r_3[x] w_2[z] r_1[z] c_2 w_1[x] c_1 w_3[x] c_3$$

\begin{enumerate}
  \item (1 point) L'une des opérations dans l'application de commerce
  électronique est la modification du prix d'un produit suite à une réduction
  de prix d'un pourcentage donné. Il faut donc lire le taux de réduction,
      le prix, et mettre à jour ce dernier. 
 Y a-t-il une ou plusieurs transactions
  de $H$ qui puissent provenir de l'exécution de cette opération? Justifiez votre
  réponse.

\begin{solution}
      Solution : Les transactions de l'exécution H sont :
      T1: r1[x] w1[x] r1[z] w1[x] c1
      T2: r2[y] r2[z] w2[z] c2
      T3: r3[z] r3[x] w3[x] c3

      La réduction du prix nécessite la lecture du taux de réduction et ensuite
      la mise-à-jour (lecture+écriture) du prix. Les transactions T2 et T3
      correspondent à cette suite d'actions. Mais les deux ne peuvent pas être
      en même temps des réductions de prix, car z ne peut pas être à la fois
      prix (pour T2) et taux de réduction (pour T3). Donc la réponse correcte
      est oui, soit T2, soit T3, mais pas les deux ensemble.
\end{solution}
  \item  (1 point) Vérifiez si $H$ est sérialisable en identifiant les conflits et en construisant le graphe de sérialisation.
\begin{solution}
      Solution : Les conflits :
      sur x : r1[x]-w3[x], w1[x]-r3[x], w1[x]-w3[x], r3[x]-w1[x]
      sur y : pas de conflit
      sur z : r3[z]-w2[z], w2[z]-r1[z]

      Le graphe de sérialisation contient un cycle T1 T3 T2 et un autre cycle T1 T3, donc H n'est pas sérialisable.
\end{solution}

  \item (1 points) Montrez que l'exécution $H$ n'évite pas les annulations en
cascade. Est-elle recouvrable? Est-il possible, en modifiant seulement la
position des Commit d'éviter les annulations en cascade?
\begin{solution}
      Solution : Après w1[x], on a plus tard r3[x] avant la fin de T1, donc H
      n'évite pas les annulations en cascade. Pareil pour w2[z] suivi de
      r1[z]. Par contre H est recouvrable, car dans les deux cas la transaction
      qui écrit est validée avant celle qui lit.

      Pour éviter les annulations en cascade, il faudrait déplacer c1 avant
      r3[x], respectivement c2 avant r1[z]. Le second déplacement est possible,
      par contre le premier est impossible, car r1[z] et w1[x] resteraient
      après c1. Donc la réponse est non.
\end{solution}

\item (1,5 points) Quelle est l'exécution obtenue par verrouillage à 
       deux phases à partir de H?

      On considère qu'il existe deux verrous pour chaque enregistrement, un de
      lecture et un d'écriture. Le relâchement des verrous d'une transaction se
      fait au Commit et à ce moment on exécute en priorité les opérations
      bloquées en attente de verrou, dans l'ordre de leur blocage.
\begin{solution}
      Solution : H : r1[x] r2[y] w1[x] r2[z] r3[z] r3[x] w2[z] r1[z] c2 w1[x] c1 w3[x] c3

      r1[x], r2[y] s'exécutent, en prenant les verrous de lecture
      w1[x] s'exécute, pas de conflit avec r1[x] (même transaction)
      r2[z] prend le verrou de lecture sur z et s'exécute
      r3[z] partage le verrou de lecture sur z avec r2[z] et s'exécute
      r3[x] bloquée par w1[x], donc T3 bloquée
      w2[z] bloquée par r3[z], donc T2 bloquée
      r1[z] partage le verrou de lecture sur z avec r2[z] et r3[z] et s'exécute
      c2 bloquée car T2 bloquée
      w1[x] a déjà le verrou et peut s'exécuter
      c1 s'exécute et relâche les verrous de T1 Þ r3[x] peut obtenir le verrou sur x et s'exécuter, par contre T2 reste bloquée
      w3[x] prend le verrou d'écriture sur x et s'exécute
      c3 s'exécute et relâche les verrous de T3 Þ toutes les opérations de T2 sont débloquées, donc w2[z] et c2 s'exécutent

      Le résultat final est donc H' : r1[x] r2[y] w1[x] r2[z] r3[z] r1[z] w1[x]
      c1 r3[x] w3[x] c3 w2[z] c2
\end{solution}

\item (1,5 points) Quelle est l'exécution obtenue par l'algorithme
       de contrôle de concurrence avec versionnement\,? Indiquez
         la (ou les)  transaction(s) rejettée(s) et expliquez
          pourquoi. Quel danger a-t-on évité?

         On suppose que chaque transaction débute
         au moment de sa première opération.

\begin{solution}
Une transaction $T$ est rejettée si elle cherche à mettre à
jour une donnée qui a été modifiée depuis le début de $T$.
C'est le cas ici pour $T_3$ qui cherche à écrire $x$.
\end{solution}
\end{enumerate}

\section{Reprise sur panne (4 points)}

\end{document}
